\documentclass[letterpaper, 10 pt, conference]{ieeeconf} 

\IEEEoverridecommandlockouts
\overrideIEEEmargins

\def\pub{false} % true for publication, false for draft
\newcommand*{\template}{../../template} 
\input{\template/preamble/preamble_conf.tex}


\title{\LARGE \bf
    Point-Wise Contraction-Based Control Based on Vector-Space Optimization under State and input Constraints
}

\author{
    Myeongseok Ryu, Kyunghwan Choi and Sesun You
    \thanks{This work was not supported by any organization}% <-this % stops a space
    \thanks{M. Ryu and K. Choi are with Cho Chun Shik Graduate School of Mobility, Korea Advanced Institute of Science and Technology (KAIST), 193 Daejeon, Republic of Korea
    {\tt\small \{dding\_98, kh.choi\}@kaist.ac.kr}}%
    \thanks{P. Misra is with the Department of Electrical Engineering, Wright State University,
    Dayton, OH 45435, USA
    {\tt\small pmisra@cs.wright.edu}}%
}


\begin{document}



\maketitle
\thispagestyle{empty}
\pagestyle{empty}


%%%%%%%%%%%%%%%%%%%%%%%%%%%%%%%%%%%%%%%%%%%%%%%%%%%%%%%%%%%%%%%%%%%%%%%%%%%%%%%%
\begin{abstract}
    Contraction theory offers 


\end{abstract}


%%%%%%%%%%%%%%%%%%%%%%%%%%%%%%%%%%%%%%%%%%%%%%%%%%%%%%%%%%%%%%%%%%%%%%%%%%%%%%%%
\section{Introduction}

\begin{itemize}
    \item Advantages of contraction-based control
    \item Challenges in existing methods
    \begin{itemize}
        \item Absence of explicit handling of constraints
        \item Requirements of prior computation of contraction metrics
        \item Limitations to online optimization in real-time applications (computational complexity)
    \end{itemize}
    \item Potentional of vector-space optimization
    \begin{itemize}
        \item Integration of constraints directly into the optimization framework
        \item Input constraints and state constraints can be handled explicitly
        \item CBF based constraints can be easily integrated
    \end{itemize}
    \item Justification of the proposed method
    \begin{itemize}
        \item Examples of effectiveness of point-wise control
        \item Examples of point-wise reference
    \end{itemize}
\end{itemize}


\section{Preliminaries}

\section{Main Results}

\begin{itemize}
\item Constraint vs Stability analysis
\end{itemize}

% state 발산?
% 추종은 괜찮

% CBF 조건 만족 -> state safe set forward invariant
% 안정성 문제 없음
% 입력같은 경우도 <- CBF+입력제한 조건 만족 -> 입력제한 set forward invariant

% contraction -> error bound 
% error bound >> e(CBF, 입력 포화량)

% - ACC 구현 기존거 돌려보기
% - 메인 리저트 쓰기 (렘마화 명확히 + CBF)

\section{Numerical Validations}

\subsection{Validation Setup}


\begin{itemize}
\item LK
\item 
\end{itemize}

\begin{equation}
    \ddtt
    \begin{pmatrix}
        \beta\\
        r\\
        \theta_f\\
        \theta_r
    \end{pmatrix}
    =
    \begin{pmatrix}
        \tfrac{1}{m v_x} \left(F_f+F_r\right) - r
        \\
        \tfrac{1}{I_z} \left(l_f F_f - l_r F_r\right)
        \\
        \tfrac{1}{\tau} \left(\omega_f - \theta_f\right)
        \\
        \tfrac{1}{\tau} \left(\omega_r - \theta_r\right)
    \end{pmatrix}
    ,
\end{equation}
where the lateral tire forces $F_f$ and $F_r$ are simplified as 
\begin{equation}
    F_f = F_{\max} \tanh\left( \tfrac{C_f}{F_{\max}} \alpha_f \right)
    ,
    F_r = F_{\max} \tanh\left( \tfrac{C_r}{F_{\max}} \alpha_r \right)
    ,
\end{equation}
with the slip angles $\alpha_f$ and $\alpha_r$ defined as
\begin{equation}
    \alpha_f(\beta,r,\theta_f) = \beta + \tfrac{l_f r}{v_x} - \theta_f
    ,
    \alpha_r(\beta,r,\theta_r) = \beta - \tfrac{l_r r}{v_x} - \theta_r
    .
\end{equation}

By denoting the state vector as $\mv{x}=\left(\beta,r,u_f,u_r\right)^\top\in\R^4$ and the control input as $\mv{u}=\left(w_f,w_r\right)\in\R^2$, the system can be represented in the SDC formulation as
\begin{equation} \label{sim:state:space}
    \ddtt \mv{x}
    =
    \mm{A}\mv{x} + \mm{B} u
    ,
\end{equation}
where $\mv{x}=\left(\beta,r,u_f,u_r\right)^\top\in\R^4$ is the state vector, 

\hfill

Following constraints are considered for the validation:
\subsubsection*{State Constraints}
To prevent tire slip, the slip angles $\alpha_f$ and $\alpha_r$ are constrained as
\begin{equation}
    \vert \alpha_f \vert \leq \alpha_{f,\max}, \quad \vert \alpha_r \vert \leq \alpha_{r,\max}
    ,
\end{equation}

\subsubsection*{Input Constraints}
Considering physical limitations of the steering system, the control inputs, steering angle rate, are constrained as
\begin{equation}
    \vert w_f \vert \leq w_{f,\max}, \quad  \vert w_r \vert \leq w_{r,\max}
    .
\end{equation}

\subsection{Controllers for Comparison}

We compare the following controllers including the proposed controller as follows:

\subsubsection*{Proposed Method}

\begin{itemize}
    \item Vector-space optimization
    \item CBF + Input saturation
\end{itemize}

\subsubsection*{CV-STEM}

\subsubsection*{Desired control law only}



\begin{table}[t]
\renewcommand{\arraystretch}{1.3}
\caption{Vehicle Dynamics Parameters.}
\centering
\begin{tabular}{c m{9.5em} c c c }
    % \begin{tabular}{c c c c c}
    \hline
    \textbf{Symbol} & \textbf{Description} & \textbf{Value} \\
    \hline
    \hline 
    $m$ & Vehicle Mass & \qty{1650}{\kilo\gram} \\
    \hline
    $(l_f,l_r)$  & Front and Rear Wheelbase & (\qty{1.11}{}, \qty{1.59}{}) \qty{}{\meter} \\  
    \hline
    $(C_f,C_r)$ & Front and Rear Cornering Stiffness & (\qty{133}{}, \qty{98.8}{}) \qty{}{\kilo\newton\per\radian} \\
    \hline
    $F_{\max}$ & Maximum Tire Force & \qty{8000}{\newton} \\
    \hline
    $I_z$ & Yaw Moment of Inertia & \qty{2315.3}{\kilo\gram\meter\squared} \\
    \hline
    $v_x$ & Longitudinal Velocity & \qty{27.7}{\meter\per\second} ($\approx$ \qty{100}{\kilo\meter\per\hour}) \\
    \hline
    \end{tabular}
    \label{table:example}
\end{table}


\subsection{Validation Results}

\section{Conclusions}



%%%%%%%%%%%%%%%%%%%%%%%%%%%%%%%%%%%%%%%%%%%%%%%%%%%%%%%%%%%%%%%%%%%%%%%%%%%%%%%%
% \section{ACKNOWLEDGMENTS}

% The authors gratefully acknowledge the contribution of National Research Organization and reviewers' comments.


%%%%%%%%%%%%%%%%%%%%%%%%%%%%%%%%%%%%%%%%%%%%%%%%%%%%%%%%%%%%%%%%%%%%%%%%%%%%%%%%

\begin{thebibliography}{99}

\bibitem{c1}
J.G.F. Francis, The QR Transformation I, {\it Comput. J.}, vol. 4, 1961, pp 265-271.

\bibitem{c2}
H. Kwakernaak and R. Sivan, {\it Modern Signals and Systems}, Prentice Hall, Englewood Cliffs, NJ; 1991.

\bibitem{c3}
D. Boley and R. Maier, "A Parallel QR Algorithm for the Non-Symmetric Eigenvalue Algorithm", {\it in Third SIAM Conference on Applied Linear Algebra}, Madison, WI, 1988, pp. A20.

\end{thebibliography}

\end{document}
